\section{Introduction}

Diffusion language models are attractive because they can update many tokens in parallel, but that same parallelism creates a structural problem: simultaneously decoded tokens are often treated as conditionally independent even when the sequence requires coordinated decisions. The resulting factorization gap is now a central bottleneck for discrete diffusion decoding. Recent work addresses the problem in several different ways, including stronger coupled output layers \citep{li2026codd}, more efficient block or semi-autoregressive decoding \citep{liu2026sdlm,wu2026fastdllmv2}, inference-time remasking or self-correction \citep{wang2025remasking,zhang2026selfcorrection}, and broader autoregressive-diffusion hybrids \citep{fathi2025unifying,sahoo2026esolm}. What remains less explored is whether the expensive dependency model should be chosen \emph{locally} rather than globally.

This paper studies that question in deliberately constrained form. Instead of claiming a new state-of-the-art diffusion language model, we isolate the controller hypothesis on a small synthetic scaffold that exposes copy-like and progression-like dependency groups. Our proposal, Adaptive Dependency Routing Diffusion (\method), routes each unresolved region at each denoising step to one of three modes: \factorized\ decoding for cheap independent updates, \coupled\ decoding for short strongly linked spans, and \micro\ decoding for tiny high-risk knots. A budget controller constrains how many tokens may receive non-factorized treatment. The design target is not ``joint modeling everywhere'' but ``pay for richer dependency handling only where the factorization gap appears largest.''

The resulting evidence is useful but mixed. Under this local scaffold, selective routing clearly improves over a fully factorized baseline. At a relaxed span-4 budget, \method\ increases token accuracy from $0.6222$ to $0.6825$ and dependency accuracy from $0.3446$ to $0.6245$ while using expensive routing for only $10.89\%$ of tokens. It also improves a longer-context stress setting. However, the strongest same-budget surrogate is not the full adaptive policy: a \coupled-only baseline slightly exceeds \method\ on token accuracy, and a random router is competitive on some dependency metrics. This means the present paper supports a feasibility claim, not a final claim that adaptive three-way routing is already the best policy.

Our contribution is therefore a controlled pilot study with three parts:
\begin{enumerate}
\item We formulate the \method\ routing problem as selective dependency allocation under explicit token and group budgets.
\item We implement a minimal synthetic evaluation scaffold that measures token accuracy, dependency recovery, mode usage, and matched-budget tradeoffs.
\item We show that selective non-factorized routing is viable, while also identifying the main failure mode: on the current scaffold, most gains can already be explained by a simpler \coupled-only policy.
\end{enumerate}

This partly negative outcome is a feature rather than an embarrassment. A clean pilot paper should tell the reader not only that selective routing can help, but also which part of the design currently carries the improvement and which part still lacks convincing evidence.

The rest of the paper is organized as follows. Section~\ref{sec:related} positions \method\ relative to recent diffusion decoding work. Section~\ref{sec:method} defines the routing mechanism and budget controller. Section~\ref{sec:setup} describes the synthetic task and matched-budget protocol. Section~\ref{sec:results} presents the frontier, ablations, and long-context stress test. Section~\ref{sec:limitations} explains why the current evidence should be read as synthetic feasibility rather than a finished dLLM result.
